% ID: SEL-D0041-OB-E04
% ESTADO: cotejada. Segunda pasada visual y borrador compilado revisados.
% Fuente: Selectividad/matematicas_aplicadas/Ua6MATEMÁTICAS APLICADAS examen.pdf
% SHA-256: b7541cc45d3e96297c29037fe17a562b7daa9cd0ee2a95c2b5e9454ad9b90f14
% Página: 2; opción B; ejercicio original 4.

\textbf{(2.5 puntos)} El peso medio de los pájaros de una determinada especie que habita en un parque natural se consideraba no inferior a 110 g, pero los biólogos del parque sostienen ahora la hipótesis de que dicho peso medio ha disminuido a consecuencia del cambio climático. Se ha tomado una muestra de 100 pájaros de esta especie y se ha obtenido un peso medio de 108 g. Se sabe que la variable que mide el peso de los pájaros de esta especie sigue una distribución Normal con desviación típica igual a 6 g.

Plantee un contraste de hipótesis $(H_0:\mu\geq110)$, con un nivel de significación del 5\%, determine la región crítica de este contraste y, utilizando ésta, razone si con ese nivel se puede aceptar que los biólogos del parque están en lo cierto.
